% !TEX program = xelatex
% 공통수학2 · Ⅰ. 도형의 방정식 · 02 직선의 위치 관계 · 2/2차시
% 교사용 지도안

\documentclass[11pt,a4paper]{article}

\usepackage{kotex}
\usepackage{iftex}
\ifPDFTeX\else
  \setmainhangulfont{Noto Serif CJK KR}
  \setsanshangulfont{Noto Sans CJK KR}
\fi

\usepackage[margin=2.1cm]{geometry}
\usepackage{parskip}
\usepackage{enumitem}
\usepackage{array}
\usepackage{tabularx}
\usepackage{booktabs}
\usepackage{xcolor}
\usepackage{amssymb}
\usepackage{tikz}
\usepackage[most]{tcolorbox}
\usepackage{fancyhdr}
\usepackage{titlesec}

\definecolor{goldc}{HTML}{B07C22}
\definecolor{goldstrong}{HTML}{8A5F16}
\definecolor{indigoc}{HTML}{2B3B57}
\definecolor{indigosoft}{HTML}{6E82A6}
\definecolor{linec}{HTML}{D6DACB}
\definecolor{surface2}{HTML}{E4E8DC}
\definecolor{notebg}{HTML}{FBF3E3}
\definecolor{inksoft}{HTML}{52604F}

\titleformat{\section}{\normalfont\sffamily\bfseries\large\color{indigoc}}{\thesection}{0.6em}{}
\titlespacing{\section}{0pt}{16pt}{8pt}

\newtcolorbox{ideabox}{enhanced, breakable, colback=white, colframe=goldc, boxrule=0.5pt, leftrule=3pt, arc=2pt, left=10pt, right=10pt, top=8pt, bottom=8pt}
\newtcolorbox{calloutbox}{enhanced, breakable, colback=white, colframe=indigosoft, boxrule=0.5pt, leftrule=3pt, arc=2pt, left=10pt, right=10pt, top=7pt, bottom=7pt}
\newtcolorbox{notebox}{enhanced, breakable, colback=notebg, colframe=goldc, boxrule=0.5pt, leftrule=3pt, arc=2pt, left=10pt, right=10pt, top=7pt, bottom=7pt}
\newtcolorbox{answerbox}{enhanced, breakable, colback=white, colframe=linec, boxrule=0.5pt, arc=2pt, left=10pt, right=10pt, top=7pt, bottom=7pt}

\newcommand{\lessonbanner}[3]{%
  \begin{tcolorbox}[enhanced, colback=indigoc, colframe=indigoc, boxrule=0pt, arc=0pt,
    left=14pt, right=14pt, top=14pt, bottom=10pt, borderline south={2.4pt}{0pt}{goldc}, fontupper=\color{white}]
    {\sffamily\footnotesize\color{white!85!black} #1}\\[4pt]
    {\sffamily\bfseries\Large #2}\\[3pt]
    {\sffamily\small\color{white!88!black} #3}
  \end{tcolorbox}
}

\pagestyle{fancy}
\fancyhf{}
\renewcommand{\headrulewidth}{0pt}
\fancyfoot[L]{\footnotesize\color{inksoft} 공통수학2 · 2026학년도 2학기 · 지평선고등학교}
\fancyfoot[R]{\footnotesize\color{inksoft} 교사용 지도안 -- 배포 금지}

\begin{document}

\lessonbanner{공통수학2 · Ⅰ. 도형의 방정식 · 1. 평면좌표와 직선의 방정식}%
  {02 직선의 위치 관계 --- 2/2차시}%
  {두 직선의 수직 조건 · 교사용 지도안 (2026학년도 2학기)}

\vspace{10pt}

\noindent
\begin{tabularx}{\linewidth}{@{}>{\bfseries\color{indigoc}}p{2.6cm}X@{}}
\toprule
대상 & 고1 · 2개 반 · 반당 13명 \\
차시 & 50분 (도입 10 · 전개 30 · 정리 10) \\
성취기준 & 두 직선의 수직 조건을 탐구하고 이해한다. \\
선행 학습 & 4차시 --- 두 직선의 평행 조건 \\
\bottomrule
\end{tabularx}

\section{학습 목표 \& Big Idea}

\begin{ideabox}
\textbf{\color{goldstrong}영속적 이해 (Big Idea)}\\
`평행하다'가 기울기의 \textit{같음}으로 번역됐다면, `수직이다'는 기울기의 \textit{곱}이라는 또 다른 대수적 조건으로 번역된다. 기하적 직관은 늘 대수식으로 바꿔 쓸 수 있다.
\vspace{4pt}

\small\color{inksoft}
핵심개념: \textbf{수직} \quad\textbullet\quad 핵심개념: \textbf{대수적 번역} \quad\textbullet\quad
일반화: \textbf{같은 도구(피타고라스 정리)가 다른 조건을 만들어낸다}
\end{ideabox}

\begin{calloutbox}
\textbf{차시 학습 목표}
\begin{itemize}[leftmargin=1.4em, itemsep=2pt, topsep=4pt]
  \item 두 직선이 수직이기 위한 조건($mm'=-1$)을 피타고라스 정리로부터 유도할 수 있다.
  \item 한 점을 지나고 주어진 직선에 수직인 직선의 방정식을 구할 수 있다.
  \item 선분의 수직이등분선의 방정식을 구할 수 있다.
\end{itemize}
\end{calloutbox}

\section{질문의 연결}

\begin{tabularx}{\linewidth}{@{}>{\bfseries\color{indigoc}\small}p{2.4cm}X@{}}
사실적 질문 & 두 직선이 수직일 때, 두 기울기 $m$, $m'$ 사이에는 어떤 관계가 성립할까? \\[6pt]
개념적 질문 & 왜 하필 곱이 $-1$일까? 피타고라스 정리는 이 조건과 어떻게 연결될까? \\[6pt]
논쟁적 질문 & 평행은 `같음', 수직은 `반대되는 곱'으로 번역된다. 관계를 이렇게 숫자 하나로 압축하는 것에는 어떤 장점과 한계가 있을까? \\
\end{tabularx}

\section{수업 흐름}

\begin{tabularx}{\linewidth}{@{}>{\centering\arraybackslash}p{1.6cm}X>{\raggedright\arraybackslash\small}p{3.1cm}@{}}
\toprule
\textbf{단계} & \textbf{활동 \& 발문} & \textbf{ATL / VTR} \\
\midrule
\textcolor{goldstrong}{\textbf{도입}}\newline\footnotesize 10분 &
\textbf{마음공부 (2분)} --- 유무념 대조.\newline
\textbf{생각 열기 (8분)} --- 세 점 O$(0,0)$, A$(1,3)$, B$(3,-1)$이 있다. 삼각형 AOB가 직각삼각형임을 확인(피타고라스 정리)하고, 직선 OA, OB의 기울기의 곱을 구해 보게 한다. &
ATL: 사고(패턴 발견)\newline VTR: See-Think-Wonder \\
\addlinespace
\textcolor{goldstrong}{\textbf{전개}}\newline\footnotesize 30분 &
\textbf{개념 형성 (14분)} --- 원점을 지나는 두 직선 $y=mx$, $y=m'x$가 수직일 때, $x=1$과의 교점 P$(1,m)$, Q$(1,m')$을 이용해 피타고라스 정리($\overline{OP}^2+\overline{OQ}^2=\overline{PQ}^2$)로 $mm'=-1$을 학생이 직접 유도.\newline
\textbf{적용 (16분)} --- 예제 2(수직인 직선의 방정식) 함께 풀이 $\to$ 문제 3 짝 활동 $\to$ `수직이등분선' 활동(두 가지 풀이 방법 --- 중점+수선 vs 두 점까지 거리가 같음 --- 비교). &
ATT: 촉진자 역할\newline ATL: 사고·비교 \\
\addlinespace
\textcolor{goldstrong}{\textbf{정리}}\newline\footnotesize 10분 &
\textbf{개념 연결 질문 (5분)} --- ``평행 조건 $m=m'$과 수직 조건 $mm'=-1$을 나란히 놓고 비교하면 무엇이 보이는가?''\newline
\textbf{성찰 (5분)} --- 감사 일기 1문장 + `직선의 위치 관계' 소단원 마무리 + 다음 차시 예고(점과 직선 사이의 거리). &
마음공부 · 감사 일기 \\
\bottomrule
\end{tabularx}

\vspace{8pt}
\begin{minipage}[t]{0.48\linewidth}
\begin{calloutbox}
\textbf{지도상의 유의점}
\begin{itemize}[leftmargin=1.2em, itemsep=2pt, topsep=4pt]
  \item 수직 조건의 유도는 반드시 원점을 지나는 평행 이동된 직선으로 단순화한 뒤 진행 --- 두 직선이 이루는 각은 평행이동에 영향받지 않음을 짚어 준다.
  \item $x$축·$y$축처럼 기울기가 정의되지 않는 경우(수직선)는 $mm'=-1$ 공식이 아니라 별도로 다뤄야 함을 안내한다.
  \item `수직이등분선' 활동은 정답을 주지 말고 두 학생(유찬·연서)의 서로 다른 풀이 방법이 같은 결과로 수렴함을 스스로 확인하게 한다.
\end{itemize}
\end{calloutbox}
\end{minipage}\hfill
\begin{minipage}[t]{0.48\linewidth}
\begin{notebox}
\textbf{인문학적 탐구 연결}\\
평행은 `나란히 가되 만나지 않음', 수직은 `정확히 반대 방향으로 만남'이다. 두 대수적 조건이 정확히 대칭적인 역수·반수 관계($m$과 $-\frac{1}{m}$)라는 사실은, 관계의 `정반대'도 결국 하나의 규칙으로 설명될 수 있다는 은유로 확장할 수 있다.
\end{notebox}
\end{minipage}

\section{형성평가 \& 답안}

\begin{answerbox}
\textbf{예제 2} \quad 점$(4,1)$을 지나고 직선 $x-2y+2=0$에 수직인 직선의 방정식\\
\textbf{\color{goldstrong}답} \; 기울기 $\frac{1}{2}$에 수직 $\to$ $m=-2$ $\to$ $y-1=-2(x-4)$ $\to$ $\boldsymbol{y=-2x+9}$
\end{answerbox}

\begin{answerbox}
\textbf{문제 3} \quad 점$(-2,2)$를 지나고 ⑴ $y=2x+5$ \quad ⑵ $2x+3y-4=0$에 각각 수직\\
\textbf{\color{goldstrong}답} \; ⑴ 기울기 $-\dfrac{1}{2}$ $\to$ $y=-\dfrac{1}{2}x+1$ \qquad ⑵ 기울기 $\dfrac{3}{2}$ $\to$ $y=\dfrac{3}{2}x+5$
\end{answerbox}

\begin{answerbox}
\textbf{활동} \quad 두 점 A$(-2,-3)$, B$(6,5)$를 잇는 선분 AB의 수직이등분선\\
\textbf{\color{goldstrong}답} \; 중점 $(2,1)$, $\overline{AB}$의 기울기 $1$ $\to$ 수직이등분선의 기울기 $-1$ $\to$ $\boldsymbol{y=-x+3}$
\end{answerbox}

\section{성취수준 (이 차시 범위)}

\begin{tabularx}{\linewidth}{@{}>{\bfseries\color{goldstrong}}p{0.9cm}X@{}}
\toprule
A & 두 직선의 수직 조건을 피타고라스 정리로부터 유도하여 설명하고, 관련된 문제를 해결할 수 있다. \\
B & 두 직선의 수직 조건을 설명하고, 관련된 문제를 해결할 수 있다. \\
C & 두 직선의 수직 조건을 이해하고, 간단한 문제를 해결할 수 있다. \\
D & 두 직선의 수직 조건을 안다. \\
E & 안내에 따라 수직인 직선의 방정식을 구할 수 있다. \\
\bottomrule
\end{tabularx}

\section{다음 차시}

\begin{calloutbox}
\textbf{03 점과 직선 사이의 거리 --- 1/2차시.} `02 직선의 위치 관계'가 여기서 마무리된다. 다음 시간부터는 점과 직선의 거리를 구하는 공식을 다루며, 오늘 배운 수직 조건이 그 공식의 유도 과정에 직접 사용됨을 도입에서 예고할 것.
\end{calloutbox}

\end{document}
